\section{Recognition and the exact ambiguity of coefficient systems}
\label{sec:coefficient-symmetries-v148}

The support inequality in the preceding section is not merely a
sufficient test in the class under consideration.  We characterize
that class intrinsically and determine exactly when its orientation
is unique.  We then compute the full linear stabilizer, including
the cases in which an orientation is lost.  These are statements
about unmarked first relations, not about a supplied matrix
presentation.

Throughout this section \(n\ge2\), \(q\ge1\), and \(D=n+q\).
Let \(E\) have dimension \(n^2\), put \(S=\Sym E\), and let
\(0\ne K\subset S_D\).  The affine cone is in \(E^*\).
Suppose that the radical of the homogeneous ideal \((K)\) is a
reduced hypersurface linearly equivalent to the determinant
hypersurface in square \(n\)-by-\(n\) matrices.  We call this the
\emph{determinant-radical condition}.  The ideal \((K)\) is taken
in the full symmetric algebra, not in its Artin truncation.
Its reduced equation gives a line \(\C\delta\subset S_n\),
and unique factorization gives the intrinsic residual subspace
\begin{equation}\label{eq:residual-recognition-v148}
                   J=K/\delta\ \subset S_q.
\end{equation}
Changing the generator \(\delta\) does not change this subspace.
The reduced singular loci of the recovered cone determine its
unordered tensor rulings.  An ordering identifies
\(E=V^*\otimes U\), with \(\dim U=\dim V=n\), up to
separate factor changes.  The subgroup induced by
\(\operatorname{GL}(U)\) on the second factor is independent of
those changes.  An ordering is called \emph{admissible} when
\(J\) is invariant under this full subgroup.  Admissibility is
therefore an intrinsic condition on an ordered ruling, and does
not ask for a chosen basis of either factor.

\begin{theorem}[Intrinsic recognition and orientation dichotomy]
\label{thm:recognition-dichotomy-v148}
Under the determinant-radical condition, an ordering is admissible
if and only if its Cauchy decomposition has
\begin{equation}\label{eq:recognized-cauchy-v148}
 J=\bigoplus_{\lambda\vdash q,\,\ell(\lambda)\le n}
         L_\lambda\otimes\mathbb S_\lambda U,
 \qquad L_\lambda\subset\mathbb S_\lambda V^*.
\end{equation}
The coefficient spaces are unique and are recovered by projection
and contraction.  Suppose at least one ordering is admissible and
put \(r_\lambda=\dim L_\lambda\) and
\(m_\lambda=\dim\mathbb S_\lambda\C^n\).  Then exactly one
of the following alternatives occurs.

\smallskip
\noindent (i) Some \(0<r_\lambda<m_\lambda\).  Precisely one
of the two orderings is admissible, and every linear isomorphism
between such relation spaces preserves that ordering.

\smallskip
\noindent (ii) Every \(r_\lambda\) is either zero or
\(m_\lambda\).  Both orderings are admissible.  The residual
space is a sum of full Cauchy summands, and transposition preserves
\(K\) after identifications of the two matrix factors.

\smallskip
Conversely, every nonzero system in
\eqref{eq:recognized-cauchy-v148} satisfies the determinant-radical
condition after multiplication by \(\delta\).  Thus the theorem
characterizes the image of the coefficient construction, without
assuming a coefficient presentation as part of the input.
\end{theorem}
\begin{proof}
The determinant polynomial is irreducible for \(n\ge2\).
Every element of \(K\) belongs to the radical \((\delta)\),
so \(K\subset\delta S_q\) and the cancellation in
\eqref{eq:residual-recognition-v148} is legitimate.  The intrinsic
recovery of the unordered factors follows from
Lemma~\ref{lem:linear-preserver-v147} and the rank-one rulings.
Changing either factor frame conjugates its full general linear
group into itself; scalar changes act as scalars on \(S_q\).
Consequently the definition of admissibility is independent of
all these choices.

For an ordering, the Cauchy formula is
\[
 S_q=\bigoplus_{\lambda\vdash q,\,\ell(\lambda)\le n}
      \mathbb S_\lambda V^*\otimes\mathbb S_\lambda U.
\]
As representations of \(\operatorname{GL}(U)\), the nonzero
\(\mathbb S_\lambda U\) are irreducible and pairwise
nonisomorphic.  Finite-dimensional rational representations of
this group over \(\C\) are completely reducible.  Hence an
invariant subspace splits across its isotypic components.  In the
\(\lambda\)-component, evaluation identifies that subspace with
\[
 \Hom_{\operatorname{GL}(U)}(\mathbb S_\lambda U,J)
                              \otimes\mathbb S_\lambda U.
\]
Schur's lemma embeds the displayed multiplicity space uniquely in
\(\mathbb S_\lambda V^*\); this is \(L_\lambda\).
This proves necessity in \eqref{eq:recognized-cauchy-v148},
including the absence of unrecorded cross-summand graphs.
Sufficiency follows since the second factor is full.  Projection
onto a Cauchy summand followed by contraction with the dual of its
second factor recovers \(L_\lambda\).

If the reverse ordering is also admissible, \(J\) is invariant
under both full factor groups.  Their product acts irreducibly on
each Cauchy summand, and these product representations are pairwise
nonisomorphic.  Complete reducibility now makes \(J\) a sum of
whole summands.  Equivalently, all \(r_\lambda\) are zero or
full.  Conversely a sum of whole summands is invariant under both
groups and under the factor exchange.  This proves the dichotomy.
One may also see uniqueness directly from the contraction ranks
\((r_\lambda,m_\lambda)\): a factor exchange reverses them,
which is impossible at a nonzero proper support.  Every linear
isomorphism of relation spaces preserves the determinant cone;
Lemma~\ref{lem:linear-preserver-v147} leaves just the two
factor alternatives.  It must therefore preserve the unique
admissible ordering in case (i).

Finally choose a nonzero \(L_\lambda\).  At an invertible
matrix \(T:U\to V\), the map \(\mathbb S_\lambda T\) is
invertible.  A nonzero covector in \(L_\lambda\) consequently
has a nonzero matrix coefficient at \(T\).  Thus the residual
coefficients do not vanish simultaneously at any invertible
matrix.  All elements of \(\delta J\) vanish on the determinant
hypersurface.  Their common zero set is exactly that hypersurface;
the Nullstellensatz and its reduced irreducible equation give
\(\rad(\delta J)=(\delta)\), as asserted.
\end{proof}

The complete reducibility and Cauchy statements here are classical
representation theory; see \cite{ChengWang}.  Their role is to
supply an intrinsic image criterion \emph{after} the matrix factors
have been recovered.  For a relation space that fails the
admissibility criterion, the theorem does not assert a one-sided
coefficient description.  In particular it is not a classification
of arbitrary homogeneous ideals or arbitrary degeneracy schemes.

\subsection{Stabilizers and the full transpose ambiguity}
Use the normal action \(T\mapsto gTh^{-1}\).  Its action on
polynomial functions is inverse pullback, so it gives an honest
representation \(\rho\) on \(E=V^*\otimes U\).
The kernel of \(\rho:\operatorname{GL}(V)\times
\operatorname{GL}(U)\to\operatorname{GL}(E)\) is the diagonal
scalar subgroup.  For a recognized system define
\[
 G_L=\{g\in\operatorname{GL}(V):
          (\mathbb S_\lambda g^{-1})^*L_\lambda=L_\lambda
                                      \text{ for every }\lambda\},
 \qquad H_K=\operatorname{Stab}_{\operatorname{GL}(E)}(K).
\]
These are closed subgroup schemes, with the stabilizer scheme
structures from their actions on Grassmannians.

\begin{theorem}[Exact coefficient stabilizer]
\label{thm:exact-coefficient-stabilizer-v148}
In case (i) of Theorem~\ref{thm:recognition-dichotomy-v148},
\begin{equation}\label{eq:proper-stabilizer-v148}
                 H_K\simeq
        (G_L\times\operatorname{GL}(U))/\mathbb G_m.
\end{equation}
In case (ii), after choosing bases for the two factors,
\begin{equation}\label{eq:full-stabilizer-v148}
 H_K\simeq
 \bigl((\operatorname{GL}_n\times\operatorname{GL}_n)
                         /\mathbb G_m\bigr)\rtimes C_2,
\end{equation}
where \(C_2\) is generated by transposition.  These are
isomorphisms of affine algebraic groups over \(\C\), and hence
of their represented functors on commutative complex algebras.

For the ungraded algebra
\(A(K)=S/((K)+\mathfrak m^{D+1})\), they give the split exact
sequence
\begin{equation}\label{eq:complete-local-aut-v148}
 1\longrightarrow U_K\longrightarrow\operatorname{Aut}(A(K))
                      \longrightarrow H_K\longrightarrow1.
\end{equation}
Here \(U_K\) is the universal higher-order substitution group
from Proposition~\ref{prop:first-relation-group-v147}, of
dimension \(n^2(\length A(K)-1-n^2)\).
\end{theorem}
\begin{proof}
Every element of \(H_K(\C)\) preserves the radical of \((K)\),
and therefore the determinant cone.  In case (i), the unique
admissible orientation excludes its transposed component.  The
remaining map is \(\rho(g,h)\).  Cancellation of the determinant
line and projection onto each Cauchy summand show that it preserves
\(J\) if and only if \(g\in G_L\); the arbitrary right map
acts on the full factor and introduces no additional condition.
The kernel is the diagonal \(\mathbb G_m\).  In case (ii),
every separate factor transformation and transposition preserves
each active full summand.  These exhaust the determinant preservers,
so they exhaust \(H_K(\C)\).  The transposition has order two
and lies outside the preserving component for \(n\ge2\), giving
the asserted semidirect product.  Its action on the left-right
quotient is conjugation by the matrix-transpose map; no particular
formula for its action on chosen factor lifts is needed.

We justify the scheme statement as well.  The homomorphism theorem
for algebraic groups identifies the quotient by the kernel with
the closed image \cite[Theorem~5.39 and Remark~5.42]{Milne2017}.
Every finite-type group scheme over a characteristic-zero field is
smooth \cite[Lemma~39.8.2]{StacksGroups}.  Thus the stabilizer and
the image are reduced closed subgroup schemes.  The equalities on
complex points just proved imply equality of their defining radical
ideals.  This gives \eqref{eq:proper-stabilizer-v148} and
\eqref{eq:full-stabilizer-v148} as algebraic-group isomorphisms,
not only as bijections on complex points.  Finally
Proposition~\ref{prop:first-relation-group-v147} applies in degree
\(D\) and gives \eqref{eq:complete-local-aut-v148} with the
stated kernel, splitting and dimension.
\end{proof}

For completeness, the omitted zero system has \(K=0\).
Then \(A(K)\) is the free truncated symmetric algebra,
\(H_K=\operatorname{GL}(E)\), and its linear automorphisms
need not preserve any matrix factors.  The determinant-radical
condition deliberately excludes this third possibility.  A full
nonzero coefficient system and a zero system therefore lose
information in different ways: the former retains unordered
matrix geometry and has exactly a transpose ambiguity, while the
latter retains no intrinsic matrix geometry at all.

\subsection{The zero/full case over a projective reduction}
The local transpose ambiguity also has an exact global form.  Fix
a nonempty set of active partitions and take the corresponding
\(\mathcal L_\lambda\) to be full, with all other coefficient
bundles zero.  Use the same construction as in
\eqref{eq:coefficient-thickening-v147}, without its strict-support
hypothesis.

\begin{proposition}[Global ambiguity for full coefficient systems]
\label{prop:full-global-ambiguity-v148}
Let \(Y_i\) be two such schemes, for the same \(n,q\) and
active partitions, over smooth connected projective bases \(B_i\)
with source bundles \(\mathcal U_i\).  They are isomorphic as
abstract complex schemes if and only if there are an isomorphism
\(\sigma:B_1\to B_2\) and an actual bundle isomorphism
\(\mathcal U_1\simeq\sigma^*\mathcal U_2\).
An isomorphism exchanging the two ruling families over a fixed
\(\sigma\) exists precisely when, in addition, the common source
bundle is projectively trivial.  Equivalently,
\[
 \mathcal U_1\simeq\mathcal M^{\oplus n},\qquad
 \sigma^*\mathcal U_2\simeq\mathcal M^{\oplus n}
\]
for a line bundle \(\mathcal M\) on \(B_1\).
Thus even the full systems recover the actual source-bundle class;
the extra ambiguity is an exchange, not an undetected line twist.
\end{proposition}
\begin{proof}
An abstract isomorphism recovers the reduction, the actual normal
bundle map, the first relation, and the unordered determinant
rulings by the proof of
Theorem~\ref{thm:general-moving-coefficients-v147}, which uses
strict support only to orient the rulings.  Over the connected
base the choice to preserve or exchange the two components is
constant.  In the preserving case, projectivity makes the left
projective map constant.  A constant lift and
Lemma~\ref{lem:actual-factor-descent-v147} give an actual
isomorphism of the source bundles.

In the exchanging case each source projective bundle is identified
with the constant left projective bundle on the other side.  It
is therefore trivial.  A vector bundle with a trivialized
projectivization is a line twist of a constant vector space:
locally lift the given projective trivialization to linear frames;
on intersections their differences are scalar matrices, whose
scalar cocycle is the required line bundle.  After pullback by
\(\sigma\), write
\(\mathcal U_i^*=\mathcal M_i^*\otimes W_i\), with \(W_i\)
constant of dimension \(n\).  The normal bundles become
\(\mathcal M_i^*\otimes(V_i\otimes W_i)\).
Their actual isomorphism induces a morphism from \(B_1\) to
the projective-linear isomorphisms between two constant
\(n^2\)-spaces.  This target is affine; the morphism is constant.
Choose its constant linear lift.  On removing it, the actual
bundle map is scalar on each projective fibre.  In local line
frames its off-diagonal entries vanish and all diagonal entries
agree.  Reducedness makes these equalities hold as bundle maps.
Their common entry is a nowhere vanishing section of
\(\mathcal Hom(\mathcal M_1^*,\mathcal M_2^*)\), and hence
\(\mathcal M_1\simeq\mathcal M_2\).  This proves necessity
of the common actual line twist, and thus of the source-bundle
isomorphism.

Conversely any actual source-bundle isomorphism, together with a
constant left identification, preserves all active full Cauchy
summands, so it gives an isomorphism of the schemes.  If the
source bundles have the common line-twist presentation above,
choose identifications of the constant factor spaces and
transpose their tensor factors, leaving the line factor intact.
This is a global actual normal map.  It preserves the determinant
line, every active full summand, and the truncation, and supplies
an exchanging isomorphism.  These constructions prove all the
claims.
\end{proof}

\subsection{The intrinsic orbit space}
Fix a rank vector \(\boldsymbol r\) with at least one proper
nonzero entry and form the product of Grassmannians
\[
 X_{\boldsymbol r}=
   \prod_{\lambda\vdash q,\,\ell(\lambda)\le n}
       \Gr(r_\lambda,\mathbb S_\lambda V^*).
\]
Zero and full factors in this product are points.  Let
\(\mathcal Z_{\boldsymbol r}\) denote the set of subspaces of
\(\Sym^D\C^{n^2}\) satisfying the determinant-radical condition
and the intrinsic admissibility test, with ranks
\(\boldsymbol r\) in the unique admissible orientation.

\begin{corollary}[Geometric orbit classification]
\label{cor:intrinsic-orbits-v148}
The coefficient construction induces bijections
\begin{equation}\label{eq:intrinsic-orbits-v148}
 X_{\boldsymbol r}(\C)/\PGL(V)
  \ \simeq\ \mathcal Z_{\boldsymbol r}/\operatorname{GL}_{n^2}(\C)
  \ \simeq\ \{A(K):K\in\mathcal Z_{\boldsymbol r}\}/\simeq.
\end{equation}
Their stabilizers, including the ineffective right factor and the
nonlinear kernel of the last passage, are precisely those in
Theorem~\ref{thm:exact-coefficient-stabilizer-v148}.
For a nonempty zero/full rank pattern there is just one relation
orbit for that pattern, with the stabilizer in
\eqref{eq:full-stabilizer-v148}.
\end{corollary}
\begin{proof}
The recognition theorem gives surjectivity of the first map.
A linear isomorphism of two such relations is orientation-preserving;
its left factor carries every \(L_\lambda\) to the corresponding
coefficient space, and its right factor is unrestricted.  Scalars
act trivially on the Grassmannians.  This proves injectivity and
the first bijection.  The first-relation orbit proposition gives
the second.  For a zero/full pattern each coefficient Grassmannian
is a point, and the full-stabilizer theorem supplies the last claim.
\end{proof}

Equation~\eqref{eq:intrinsic-orbits-v148} is a statement about
geometric orbit sets with the explicitly computed affine
stabilizers.  It does not identify nonreduced moduli stacks or
assert that every infinitesimal deformation remains in this class.
For smooth projective reduced bases, its uniquely oriented case
combines with actual factor descent in
Theorem~\ref{thm:general-moving-coefficients-v147}.  The following
pencil-native application shows that this global step separates an
entire covering problem that fibrewise classification cannot see.
